\documentclass[10pt,a4paper]{amsart}
\usepackage[margin=19mm]{geometry}
\usepackage{amsmath,amssymb,mathtools}
\usepackage[scr=rsfs,cal=euler]{mathalfa}
\usepackage{xfrac}
\usepackage[unicode,hidelinks,bookmarks=false]{hyperref}
\newcommand{\ie}{i.e.\ }
\newcommand{\cf}{cf.\ }
\newcommand{\resp}{respectively\ }
\newcommand{\Subsection}{\subsection}
\providecommand{\nobreakdash}{\nobreak\hbox{-}}
\setlength{\emergencystretch}{2em}
\multlinegap0pt
\usepackage{bm}
\usepackage{bbm}
\usepackage{dsfont}
\usepackage{xspace}
\usepackage{cancel}
\usepackage{mathtools}
\usepackage{amsthm}
\makeatletter
\@ifundefined{theorem}{\newtheorem{theorem}{Theorem}}{}
\makeatother
\newcommand {\Q}{{\mathbb{Q}}}
\newcommand {\Z}{{\mathbb{Z}}} 
\newcommand{\calA}{\mathcal {A}}
\newcommand{\calF}{\mathcal {F}}
\newcommand{\calG}{\mathcal {G}}
\newcommand{\calR}{\mathcal {R}}
\newcommand{\calX}{\mathcal {X}}
\title[Number of integers represented by families of binary forms I]{Number of integers represented by families of binary forms III: fewnomials}
\author{Etienne Fouvry, Michel Waldschmidt}
\date{Source: arXiv:2509.08335v1 (CC BY 4.0)}
\begin{document}
\maketitle

\noindent\emph{Source and licence.} Number of integers represented by families of binary forms III: fewnomials. Version arXiv:2509.08335v1 (CC BY 4.0), distributed under the Creative Commons Attribution 4.0 International licence, as recorded in the arXiv licence metadata for this exact version and shown on the abstract page https://arxiv.org/abs/2509.08335v1. Full rights notice: licenses/arxiv-2509.08335-CC-BY-4.0.txt.\par\smallskip

\noindent\textit{Editorial scope.} Counted unit: the Theorem whose source label is \texttt{Th:majorationasymptotique} in the article (source0.tex lines 741--767, source label \texttt{Th:majorationasymptotique}), delivered together with the article's own complete original proof of it (source0.tex lines 770--843, one \texttt{proof} environment of 74 lines). No mathematics is written, repaired or re-derived: statement and proof are the article's own text line for line. Required context retained uncounted: Equation:New1 (lines 407--410); Th:principal (lines 474--504); Equation:diophantienne (lines 730--732). Every definition, display and label of those lines is retained unchanged. Omitted from this package and not redistributed: 1--406, 411--473, 505--729, 733--740, 768--769, 844--3095. Nothing inside a retained excerpt is omitted or altered: every retained body is delivered exactly as the article's lines are written, so no omission receipt of any kind accompanies this package. Citations: the retained excerpts carry no citation command at all, so no bibliography entry of the article is transcribed and no reference is redistributed. Reference adaptation: none was needed and none was made; every reference inside the delivered unit resolves inside it, so no unresolved reference or citation is produced. Printed slips kept literal: no printed word, symbol, hypothesis or number of the counted statement or of its proof was altered. Layout and class: the article's own class is \texttt{article} with packages including inputenc, amsmath, amsthm, amssymb, enumitem, graph icx, fontenc, hyperref, xcolor; the wrapper is the lane's pinned \texttt{amsart} editorial wrapper at 10pt on A4, which restates the theorem-like environments the retained text needs and adds the article's own macro definitions that the retained text needs (\texttt{\textbackslash Q}, \texttt{\textbackslash Z}, \texttt{\textbackslash calA}, \texttt{\textbackslash calF}, \texttt{\textbackslash calG}, \texttt{\textbackslash calR}, \texttt{\textbackslash calX}), with extra wrapper packages (bm, bbm, dsfont, xspace, cancel, mathtools). The delivered numbering is the unit's own. Assets: no retained span contains a figure environment, a graphics-inclusion command, a PDF-page inclusion, a table or a TikZ picture; no image file is copied into this package, the package declares no asset dependency and no image file is redistributed with it. Licence: version arXiv:2509.08335v1 of this article is distributed under the Creative Commons Attribution 4.0 International licence, as recorded in the arXiv licence metadata for this exact version and shown on the abstract page https://arxiv.org/abs/2509.08335v1. A full rights notice is delivered at licenses/arxiv-2509.08335-CC-BY-4.0.txt. Characters: every retained excerpt is pure ASCII, so the delivered engine, XeLaTeX with its default Latin Modern text font, reports no missing character.
 \begin{equation}\label{Equation:New1}
 \max_{F\in \mathcal F_d}
\left(2\, \calA ^\star(F)+1\right)^{r+1}.
\end{equation}

 \begin{theorem} \label{Th:principal}
Let $\epsilon>0$ and let $r\geqslant 2$ be an integer. Let $\mathcal F= \mathcal F_{\boldsymbol{ \mathcal D}}$ be the family of binary fewnomials attached to the system 
of data ${\boldsymbol{ \mathcal D}} =(r, (\mathcal E_k))$. 
There exists a constant $\eta>0$ which depends only on $\epsilon$ and $r$ with the following property. 
Assume that there exists $d_0\geqslant 3$ such that, for all $d\geqslant d_0$,
\begin{equation}\label{Equation:majorationCalAeta}
\max_{F\in\calF_d}\calA^\star(F)\leqslant \exp(\eta d/\log d).
\end{equation}
Then 
 
(a) For all $m\in\Z\smallsetminus\{-1,0,1\}$ and all $d\geqslant 3$ which is a multiple of $r$, the set $\calG_{\geqslant d}(m)$ is finite. Moreover, for all $d\geqslant d_0$ which is a multiple of $r$ and all $\epsilon>0$, there exists a constant $c$ depending only on $r,d,\epsilon$ such that, for $|m|\geqslant 2$,
\[
\sharp \calG_{\geqslant d}(m)\leqslant c |m|^{(1/d)+\epsilon }.
\]

(b) 
Assume that $\mathcal F$ is a $\Q$--homography--free set. Then 
for all $d\geqslant 3$ which is a multiple of $r$, we have, for $N\to\infty$,
\begin{equation}\label{Equation:calRdN}
\sharp \calR_{\geqslant d}(N)=
\left(\sum_{F\in\calF_d} C_F\right) N^{2/d}+ O_{\epsilon,r,d} \Bigl(N^{\max\{\vartheta_{d}+\epsilon,2/d^\dag\}} \Bigr).
\end{equation}

(c) 
Assume that $\mathcal F$ is a $\Q$--homography--free set of $\Q$--rigid forms. Then 
for all $d\geqslant 3$ which is a multiple of $r$, we have, for $N\to\infty$,
\[
\sharp \calR_{\geqslant d}(N)= \frac 1{(d,2)}
\left(\sum_{F\in\calF_d} A^\star_F\right) N^{2/d}+ O_{\epsilon,r,d} \Bigl(N^{\max\{\vartheta_{d}+\epsilon,2/d^\dag\}} \Bigr).
\]
\end{theorem}

\begin{equation}\label{Equation:diophantienne}
 |F(x,y)|\geqslant \calX^{d - C r^{4r} (\log d) (\log \calA^\star(F))}.
\end{equation}

 \begin{theorem}\label{Th:majorationasymptotique} 
Under the assumptions of Theorem \ref{Th:principal},
let $\lambda$ and $\mu$ be two real numbers satisfying $\lambda>2$ and
\begin{equation}\label{Equation:mu}
 0< \mu<\frac{\lambda-2}{Cr^{4r}\lambda}\cdotp
\end{equation}
Let $d_0\geqslant 3$ be an integer. Assume that the condition 
 \begin{equation}\label{Equation:majorationCalAmu}
\calA^\star(F)\leqslant \exp(\mu d/\log d)
\end{equation}
 is satisfied for all $d\geqslant d_0$ and all $F\in\calF_d$. Then 
 
 \vskip .3cm
 \noindent
(a) For all $m\in\Z\smallsetminus\{-1,0,1\}$ and all $d\geqslant 3$ multiple of $r$, the set $\calG_{\geqslant d}(m)$ is finite. 
Furthermore, for all $(\lambda, \mu,r,d)$ as above with $d\geqslant d_0$, there exists a constant $c_1$, 
only depending on $( \lambda, \mu, r,d)$, such that, for every $\vert m\vert \geqslant 2$, one has the inequality
\[
\sharp \mathcal G_{\geqslant d} (m) \leqslant c_1 \vert m\vert^{\lambda /(2d)}.
\]
 
 \vskip .3cm
 \noindent (b) For all $(\lambda, \mu,r,d)$ as above with $d\geqslant d_0$, there exists $c_2$ depending only on $(\lambda, \mu, r, d)$, such that, for all $ N \geqslant 2$, one has the inequality 
 \[
 \sharp \mathcal R_{\geqslant d} (N) \leqslant c_2N^{\lambda /d}.
 \]
 \end{theorem}

\begin{proof} Define $\lambda'$ by the equality 
\[
\mu = \frac{\lambda' -2} {C r^{4r} \lambda'}\cdotp
\]
By \eqref{Equation:mu}, we have the equalities $2 < \lambda' < \lambda$. The number $\theta := \lambda'/2$ satisfies
\[
\theta >1 \text{ and } 1-\frac 1 \theta = \frac{\lambda'-2}{\lambda'} = Cr^{4r}\mu.
\]
Let $d\geqslant 3$ be a multiple of $r$, written as $d=kr$. Let $m\in \Z$, $\vert m \vert \geqslant 2$ be such that $\mathcal G_{\geqslant d} (m) \neq \emptyset$: there exists $(x,y, F) \in {\mathcal G}_{\geqslant d} (m)$ such that $m=F(x,y)$, where $F\in \mathcal F_{k'r} $ (with $k'\geqslant k$) is the binary form
\[
a_0 X^{k'r} +a_1 X^{k'(r-1)} Y^{k'}+\cdots +a_{r-1} X^{k'} Y^{k'(r-1)} + a_r Y^{k'r}.
\]
Let $\mathcal X = \max \{ \vert x \vert, \, \vert y \vert \}$. By hypothesis, we have $\mathcal X \geqslant 2$. The lower bound \eqref{Equation:diophantienne} yields 
\[
\vert m \vert \geqslant \mathcal X^{k'r -Cr^{4r} (\log (k'r))(\log \mathcal A^\star (F))}.
\]
Assume now $d\geqslant d_0$. 
From \eqref{Equation:majorationCalAmu} we deduce the upper bound
\[
Cr^{4r} (\log (k'r)) (\log \mathcal A^\star (F)) \leqslant Cr^{4r+1} \mu k' = k'r \left( 1 -\frac 1 \theta \right),
\]
hence
\[
\mathcal X^{k'r} \leqslant \vert m \vert^\theta.
\]
Define
\[
M_0:= \left\lfloor \frac{\theta \log \vert m \vert}{r \log 2}\right\rfloor.
 \] 
Thanks to the inequalities $\mathcal X \geqslant 2$ and $k'\geqslant k$ we deduce $M_0\geqslant k$ and
\begin{equation}\label{Equation:M0}
k'\leqslant M_0 \qquad \text{ and } \qquad \mathcal X \leqslant \vert m \vert ^{\theta/ (k'r)}.
\end{equation}
Given $x$ and $m$, the number of $y^{k'}$ such that $F(x,y)=m$ is at most $r$, hence the number of such $y$ is at most $2r$. This shows that the set $\mathcal G_{\geqslant d} (m)$ is finite for all $m$ with $|m|\geqslant 2$ and $d\geqslant d_0$. Since, for all $d \geqslant 3$ all the sets $\mathcal F_d$ are finite, we deduce from Thue's Theorem on the finiteness of the number of solutions of Thue's equation that the set $\mathcal G_{\geqslant d} (m)$
is finite for any $d\geqslant 3$.

We also deduce for $d\geqslant d_0$ that the number of elements in $\mathcal G_{\geqslant d} (m)$ is bounded by
\begin{equation}\label{Equation:sharpGd(m)}
\sharp \, \mathcal G_{\geqslant d} (m) \leqslant 4r \vert m \vert^{\theta/d} \sum_{k'=k}^{M_0} \sharp \mathcal E_{k'}. 
\end{equation}
Using the inequality \eqref{Equation:New1} under the form 
\[
\sharp \mathcal E_{k'} \leqslant 3^{r+1} \max_{F\in \mathcal F_{k'r}} (\mathcal A^\star(F))^{r+1},
\]
together with the hypothesis (2.3) we deduce 
\[
\sharp \mathcal E_{k'} \leqslant 3^{r+1}\, \exp \left(\mu k' r(r+1) /(\log (k'r))\right). 
\]
Combining with \eqref{Equation:sharpGd(m)} and the upper bound $\theta<\lambda/2$, we deduce 
the inequality
\begin{align*}
\sharp \, \mathcal G_{\geqslant d} (m) &\leqslant 4r \vert m \vert^{\theta/d} \cdot 
M_0 \cdot 3^{r+1}\, \exp \left(\mu M_0 r(r+1) /(\log (M_0r))\right) \\
& \leqslant c_1 \vert m\vert^{\lambda/(2d)}.
\end{align*} 
This completes the proof of the item (a).

The proof of the item (b) has similarities and works as follows. Write $d=kr$ and let $m$ be an element of $\mathcal R_{\geqslant d}(N)$.
It satisfies $\vert m \vert \leqslant N$ and it can be written as 
$m= F(x,y)$ for some $(x,y)$ such that $ \mathcal X \geqslant 2$, for some $k'\geqslant k$ and some $F\in \mathcal F_{k'r}$.
However the inequalities \eqref{Equation:M0} imply
\[
k'\leqslant M_1\quad \text{ and } \quad \max\{|x|,|y|\} \leqslant N^{\theta/(kr)}
\quad \text{ where } \quad 
M_1:=
\left\lfloor
 \frac{\theta \log N}{r\log 2} \right\rfloor.
\]
Thus we have
\[
\sharp \mathcal R_{\geqslant d} (N)\leqslant \left(1+2 N^{\theta/(kr)}\right)^2 \sum_{k'=k}^{M_1} \sharp \mathcal E_{k'},
\]
and the item (b) follows.
\end{proof}

\end{document}
